% ==========================================
% 3. State of the Art
% ==========================================
%
% Status: living section. The structure (themes) is fixed early, from the
% research questions; the content grows during the whole project; the synthesis
% and positioning are written last.
%
% Working method: each author owns one or two themes, takes notes in
% notebook/Literature_Review.md, then turns them into LaTeX. One person performs
% the integration pass so that the section reads as a single text.
%
% Rules:
%   - We organize by theme or approach, not paper by paper.
%   - We cite only sources that we have actually read (at least abstract,
%     method and conclusion).
%   - Every source is added to references.bib as soon as it is read.

\section{State of the Art}
\label{sec:state-of-the-art}

% ------------------------------------------
% Short framing paragraph
% ------------------------------------------
% TODO [NOW]: We announce the themes covered and explain how they relate to the
%             research questions of Section~\ref{sec:rq-hypotheses}.

\subsection{Literature Search Method}
\label{sec:sota-method}
% TODO [DURING]: We state how the literature was found: databases and search
%                engines used (e.g. Google Scholar, IEEE Xplore, ACM DL, arXiv),
%                keywords, period covered, and the inclusion and exclusion
%                criteria.
% TODO [CHECK]: We state the limits of the search (language, access to papers),
%               so that the review is reproducible.

\subsection{Theme A: [to be defined]}
\label{sec:sota-theme-a}
% TODO [NOW]: We choose the themes from the research questions (not from the
%             papers we happen to know) and rename these subsections.
% TODO [DURING]: For each theme, we present the main approaches, their principle,
%                their assumptions, their strengths and their weaknesses.
% TODO [DURING]: We compare the works within a theme on common criteria (data,
%                performance, limits) rather than summarizing them one by one.
% TODO [DURING]: We cite foundational works as well as recent ones, and prefer
%                primary sources to secondary ones.

\subsection{Theme B: [to be defined]}
\label{sec:sota-theme-b}
% TODO [DURING]: Same structure as Theme A.

\subsection{Theme C: [to be defined]}
\label{sec:sota-theme-c}
% TODO [DURING]: Same structure as Theme A. We add or remove themes as needed.

\subsection{Synthesis and Positioning}
\label{sec:sota-synthesis}
% TODO [LAST]: We summarize the compared approaches in one table (rows: works or
%              approaches, columns: comparison criteria).
% TODO [LAST]: We state the gaps explicitly: what existing work does not solve,
%              each gap linked to a research question.
% TODO [LAST]: We position our study: what it reuses, what it changes, what it
%              adds. We make claims proportionate to the evidence (no
%              unsupported "novel" or "first").
% TODO [CHECK]: The gaps listed here match the problem statement
%               (Section~\ref{sec:problem-statement}) and the introduction.
% TODO [CHECK]: Every statement about existing work carries a citation, and every
%               cited key exists in references.bib.
